\subsection{PARACOMPACT SPACES}

We recall (Chapter I, § 9, no. 10) that a topological space X is said to be paracompact if it is Hausdorff and if every open covering of X has a locally finite open refinement.

PROPOSITION 4. Every paracompact space is normal.

This is a consequence of the following lemma:

LEMMA 2. Let A, B be two disjoint closed subsets of a paracompact space X. If for each \(x \in \mathbf{A}\) there is an open neighbourhood \(\mathbf{V}_x\) of \(x\) and an open neighbourhood \(\mathbf{W}_x\) of B which do not intersect, then there exists an open neighbourhood T of A and an open neighbourhood U of B which do not intersect.

Assuming the truth of this lemma for the moment, we can apply it to the case where B consists of a single point, because X is Hausdorff, and it shows then that X is regular. We can then apply Lemma 2 again to any two disjoint closed subsets of X, and this shows that Axiom ( \(O_{V}^{\prime}\) ) is satisfied.

To prove the lemma, consider the open covering of X consisting of \(\mathbb{C}\mathbf{A}\) and the \(\mathbf{V}_x\) , where \(x \in \mathbf{A}\) ; let \((\mathbf{T}_t)_{t \in \mathbf{I}}\) be a locally finite open refinement of this covering. Then, by definition, if \(\mathbf{A} \cap \mathbf{T}_t \neq \emptyset\) there exists \(x_t \in \mathbf{A}\) such that \(\mathbf{T}_t \subset \mathbf{V}_{x_t}\) . Let T be the open set which is the union of the \(\mathbf{T}_t\) which meet A, and let us show that there is an open neighbourhood U of B which does not meet T. For each \(y \in \mathbf{B}\) there is an open neighbourhood \(\mathbf{S}_y\) of \(y\) which meets only a finite number of sets \(\mathbf{T}_t\) ; let J be the finite subset of I consisting of those indices t such that \(\mathbf{T}_t\) meets both \(\mathbf{S}_y\) and A; if we put \(\mathbf{U}_y = \mathbf{S}_y \cap \bigcap_{t \in \mathbf{J}} \mathbf{W}_{x_t}\) , then \(\mathbf{U}_y\) is an open neighbourhood of \(y\) which meets none of the \(\mathbf{T}_t\) which meet A, and hence \(\mathbf{U}_y \cap \mathbf{T} = \emptyset\) . Let \(\mathbf{U} = \bigcup_{y \in \mathbf{B}} \mathbf{U}_y\) ; then \(\mathbf{U}'\) is an open neighbourhood of B which does not meet T, and the lemma is proved.

There exist normal spaces which are not paracompact (Exercise 19).

COROLLARY 1. Given any open covering \((\mathbf{A}_t)_{t \in \mathbf{I}}\) of a paracompact space \(\mathbf{X}\) , there exists a continuous partition of unity \((f_t)_{t \in \mathbf{I}}\) on \(x\) , subordinate to the covering \((\mathbf{A}_t)\) .

Let \((\mathbf{U}_{\lambda})_{\lambda \in \mathbf{L}}\) be a locally finite open covering of \(\mathbf{X}\) which refines the covering \((\mathbf{A}_{\iota})_{\iota \in \mathbf{I}}\) ; then there is a mapping \(\varphi : \mathbf{L} \to \mathbf{I}\) such that \(\mathbf{U}_{\lambda} \subset \mathbf{A}_{\varphi(\lambda)}\) for each \(\lambda \in \mathbf{L}\) . By Propositions 3 and 4, there exists a continuous partition of unity \((g_{\lambda})_{\lambda \in \mathbf{L}}\) subordinate to \((\mathbf{U}_{\lambda})\) . For each \(\iota \in \mathbf{I}\) , put

\[
f _ {t} = \sum_ {\varphi (\lambda) = t} g _ {\lambda};
\]

this sum is defined and continuous since the supports of the \(g_{\lambda}\) form a locally finite covering; moreover, the union \(\mathbf{B}_{t}\) of the supports of the \(g_{\lambda}\) such that \(\varphi(\lambda)=\iota\) is closed, by Proposition 4 of Chapter I, § 1, no. 6, and is contained in \(\mathbf{A}_{t}\) . Since we have \(f_{t}(x)=0\) whenever \(x\in\complement\mathbf{B}_{t}\) , the support of \(f_{t}\) is contained in \(\mathbf{B}_{t}\) , and therefore in \(\mathbf{A}_{t}\) . On the other hand, the family \(\mathbf{B}_{t}\) is locally finite, because for each \(x\in\mathbf{X}\) there is a neighbourhood V of \(x\) and a finite subset H of L such that V∩U \(_{\lambda}=\emptyset\) for all \(\lambda\notin\mathbf{H}\) , and it follows therefore that V∩B \(_{t}=\emptyset\) for all \(\iota\notin\varphi(\mathbf{H})\) . Finally, for each \(x\in\mathbf{X}\) we have

\[
\mathrm{I} = \sum_ {\lambda \in \mathbf {L}} g _ {\lambda} (x) = \sum_ {t \in \mathbf {I}} \left(\sum_ {\varphi (\lambda) = t} g _ {\lambda} (x)\right) = \sum_ {t \in \mathbf {I}} f _ {t} (x),
\]

and the proof is complete.

COROLLARY 2. If F is a closed subset of a paracompact space X, then every neighbourhood of F in X contains a closed (and therefore paracompact) neighbourhood of F.

By Proposition 16 of Chapter I, § 9, no. 10, any closed subspace of X is paracompact; the corollary therefore follows from Proposition 4 and Axiom (O \(_{V}\) '').
% semantic-fragments: legacy-input-seam
\relax{}
